% Conteudo para Overleaf:
% Inclua no seu main.tex com: % Conteudo para Overleaf:
% Inclua no seu main.tex com: % Conteudo para Overleaf:
% Inclua no seu main.tex com: % Conteudo para Overleaf:
% Inclua no seu main.tex com: \input{docs/estudo.tex}
%
% Obs.: texto em formato LaTeX (secoes e subsecoes) para estudo teorico.

\section{Estudo - Tecnologias e Conceitos do Projeto}

\subsection{Visao Geral}
Este projeto implementa um \textbf{sistema de deteccao de quedas} baseado em aprendizado profundo e visao computacional, com integracao de notificacoes via \textbf{ESP32} (hardware) e \textbf{aplicativo mobile} (React Native/Expo). A arquitetura ponta a ponta envolve:
captura de video por camara, processamento em janelas temporais, inferencia com uma rede \textbf{CNN+LSTM}, e distribuicao de alertas usando \textbf{MQTT} (com broker Mosquitto) para acionamento local (buzzer/LEDs) e remoto (celular).

\subsection{Objetivos do Estudo}
Compreender e relacionar:
\begin{itemize}
  \item Tecnicas de \textbf{deep learning} para classificacao de eventos em video.
  \item Processamento de video com \textbf{OpenCV} e preparo de dados (ETL).
  \item Padroes de rede neural: \textbf{MobileNetV2}, \textbf{TimeDistributed}, \textbf{LSTM} e classificacao binaria.
  \item Treinamento supervisionado: hiperparametros, callbacks e regularizacao.
  \item Avaliacao estatistica: matriz de confusao, precision, recall e F1-score.
  \item Inferencia e otimizacao para \textbf{CPU} com \textbf{TFLite} e quantizacao.
  \item Integracao IoT: \textbf{ESP32}, comunicacao \textbf{Serial/UART}, \textbf{MQTT} e mensageria em JSON.
  \item App mobile: \textbf{React Native}, \textbf{Expo}, WebSocket para MQTT e persistencia local com AsyncStorage.
  \item Reprodutibilidade: logs de treinamento e controle de artefatos.
\end{itemize}

\subsection{Stack Tecnologico (Resumo)}
\begin{itemize}
  \item \textbf{Backend IA e Visao Computacional:} Python, TensorFlow/Keras, OpenCV, NumPy, scikit-learn.
  \item \textbf{Treinamento e Avaliacao:} tf.data, callbacks (EarlyStopping, ReduceLROnPlateau, ModelCheckpoint), matplotlib.
  \item \textbf{Inferencia Otimizada:} TFLite com quantizacao INT8 (quando disponivel).
  \item \textbf{IoT:} ESP32 (Arduino), Serial (UART) via pyserial e MQTT via paho-mqtt.
  \item \textbf{Broker:} Eclipse Mosquitto com listener MQTT (1883) e WebSocket (9001) para o app.
  \item \textbf{App Mobile:} React Native + Expo + TypeScript.
  \item \textbf{Persistencia e Media:} AsyncStorage (historico/config), audio em loop e vibra (haptics) via bibliotecas do Expo.
\end{itemize}

\section{Tema Central: Deteccao de Quedas em Video}

\subsection{Por que Video e Tempo Importam}
Quedas sao eventos dinamicos: nao se resumem a uma postura unica. Por isso, e necessario modelar a \textbf{evolucao temporal} do movimento. Em vez de classificar frames isolados, o sistema trabalha com \textbf{sequencias} (janelas temporais) de frames.

\subsection{Abordagem de Classificacao Binaria}
O problema e formulado como classificacao binaria:
\begin{itemize}
  \item \textbf{Normal} (0)
  \item \textbf{Fall} (1)
\end{itemize}
A rede produz um valor no intervalo $[0,1]$ via ativacao \textbf{sigmoid}, interpretado como probabilidade da classe Fall.

\section{Modelos de Deep Learning: CNN+LSTM}

\subsection{Conceito de CNN (Convolutional Neural Network)}
CNNs extraem caracteristicas espaciais relevantes (bordas, texturas e padrões visuais). Em classificacao de imagens, uma CNN mapeia uma entrada de dimensoes $H \times W \times C$ para uma representacao vetorial.

\subsection{Transfer Learning com MobileNetV2}
O projeto utiliza \textbf{MobileNetV2} como extrator de caracteristicas por:
\begin{itemize}
  \item eficiencia computacional (arquiteturas de mobile para menor latencia);
  \item uso de pesos pre-treinados no ImageNet (transfer learning).
\end{itemize}
Transfer learning significa congelar a base inicialmente (ou em parte) para reduzir tempo e necessidade de dados rotulados.

\subsection{TimeDistributed e Processamento de Sequencias}
Para cada sequencia de $T=20$ frames, a CNN e aplicada individualmente a cada frame. A camada \textbf{TimeDistributed} replica a operacao da CNN ao longo do eixo temporal:
\begin{itemize}
  \item entrada: $(T, H, W, C)$ por amostra;
  \item saida: vetores de caracteristicas por frame (para alimentacao temporal).
\end{itemize}

\subsection{LSTM (Long Short-Term Memory)}
A LSTM modela dependencia temporal e aprendizados de como caracteristicas mudam ao longo do tempo. O componente LSTM tende a lidar com informacoes de longo prazo atraves de \textbf{portas de controle} (input/forget/output) e estados internos.

\subsection{Regularizacao (Dropout)}
O modelo aplica \textbf{Dropout} na camada de saida da LSTM para reduzir overfitting. Em geral, Dropout desativa aleatoriamente parte das unidades durante o treinamento, forçando o modelo a aprender representacoes mais robustas.

\subsection{Funcao de Perda e Otimizador}
O projeto utiliza:
\begin{itemize}
  \item \textbf{Loss:} binary crossentropy (adequada a classificacao binaria);
  \item \textbf{Otimizador:} Adam com learning rate configurado.
\end{itemize}

\section{Preparacao de Dados (ETL) para Video}

\subsection{Dataset UR Fall Detection}
O dataset usado contem:
\begin{itemize}
  \item sequencias de quedas simuladas (Fall);
  \item atividades diarias (ADL/Normal).
\end{itemize}
O sistema assume que videos podem ser organizados por pastas/classe.

\subsection{Pipeline de Pre-processamento}
O fluxo de preparo de dados envolve:
\begin{itemize}
  \item identificar arquivos por busca recursiva;
  \item extrair frames ou converter sequencias para um formato de video padronizado;
  \item redimensionar imagens para o tamanho exigido pelo modelo;
  \item normalizar valores (ex.: escala para $[0,1]$);
  \item construir amostras via \textbf{sliding window} (janelas temporais sobre o video).
\end{itemize}

\subsection{Sliding Window e Passo Temporal}
O sliding window cria varias amostras a partir de um unico video. O passo temporal (stride) controla a sobreposicao entre janelas:
\begin{itemize}
  \item stride menor => mais amostras, maior custo computacional;
  \item stride maior => menos amostras, potencialmente menos cobertura temporal.
\end{itemize}

\subsection{Data Augmentation}
Para melhorar generalizacao, o projeto aplica aumentos aleatorios em tempo real em sequencias, como:
\begin{itemize}
  \item flip horizontal;
  \item variacao de brilho.
\end{itemize}
Essa etapa reduz overfitting e melhora robustez para variacoes do ambiente.

\subsection{Split Estratificado (Treino/Teste)}
A separacao treino/teste e estratificada para preservar proporcao de classes em cada subconjunto. Isso evita que o modelo seja treinado com uma distribuicao distorcida.

\section{Treinamento do Modelo}

\subsection{Hiperparametros}
Aspectos estudaveis incluem:
\begin{itemize}
  \item numero de epocas;
  \item batch size;
  \item learning rate;
  \item tamanho da sequencia (T=20);
  \item dimensao de entrada (H=W=128).
\end{itemize}

\subsection{tf.data e Eficiência}
O uso de tf.data ajuda a:
\begin{itemize}
  \item prefetch e pipelining de dados;
  \item reduz overhead de carregamento;
  \item suportar augmentation em tempo real.
\end{itemize}

\subsection{Callbacks}
Callbacks implementados:
\begin{itemize}
  \item \textbf{ModelCheckpoint:} salva o melhor modelo conforme val\_accuracy;
  \item \textbf{EarlyStopping:} interrompe treinamento quando val\_loss para de melhorar;
  \item \textbf{ReduceLROnPlateau:} reduz learning rate quando val\_loss estagna.
\end{itemize}

\subsection{Mixed Precision (Opcional)}
Quando habilitado, mixed precision acelera treinamento em hardware compatível, reduzindo memoria e aumentando throughput. No projeto, o modo esta configurado como desativado por padrao.

\section{Avaliacao: Metrics e Analise de Erros}

\subsection{Accuracy e Loss}
Accuracy mede a proporcao de amostras corretamente classificadas. Loss quantifica o erro durante treinamento e avaliacao, sendo util para observar convergencia.

\subsection{Precision, Recall e F1-score}
Precision (para a classe prevista) e Recall (para a classe real) sao essenciais em cenarios desbalanceados. F1-score combina ambos em uma unica medida.

\subsection{Matriz de Confusao}
A matriz de confusao $C$ resume:
\begin{itemize}
  \item True Positives (acertos de Fall);
  \item True Negatives (acertos de Normal);
  \item False Positives (Normal previsto como Fall);
  \item False Negatives (Fall previsto como Normal).
\end{itemize}
Ela e a principal fonte visual de entendimento de erros.

\subsection{Leitura dos Graficos}
Os graficos de curves.png mostram:
\begin{itemize}
  \item train loss vs val loss;
  \item train accuracy vs val accuracy.
\end{itemize}
Quando val e train evoluem de forma semelhante, tende a indicar boa generalizacao (sem overfitting significativo).

\section{Inferencia em Tempo Real e Otimizacao}

\subsection{Buffer Circular de Sequencias}
Durante inferencia, o sistema acumula frames em um buffer ate completar T=20. A cada passo de inferencia, a janela atual e processada.

\subsection{Skip Frames}
O uso de inferencia a cada N frames (skip frames) reduz custo computacional e viabiliza execucao em CPU.

\subsection{Limiar e Confirmacao Temporal}
Para reduzir falsos positivos temporais, aplica-se:
\begin{itemize}
  \item um limiar mais conservador para considerar Fall;
  \item confirmacao temporal: so dispara Fall apos N predições consecutivas.
\end{itemize}

\subsection{TFLite e Quantizacao INT8}
O modelo e convertido para TFLite para desempenho. Quantizacao INT8 reduz tamanho e acelera em CPU, com possivel impacto na precisao (avaliar sempre).

\section{Visao Computacional com OpenCV}

\subsection{Captura de Video}
OpenCV utiliza VideoCapture para:
\begin{itemize}
  \item acessar webcam (indice 0);
  \item acessar arquivo de video (caminho do arquivo).
\end{itemize}

\subsection{Redimensionamento e Normalizacao}
A entrada do modelo precisa ser padronizada:
\begin{itemize}
  \item resize para HxW;
  \item conversao para float32 e normalizacao para [0,1];
  \item empilhamento temporal em formato correto para a rede.
\end{itemize}

\subsection{Visualizacao}
O sistema usa overlays com texto e paines para:
\begin{itemize}
  \item visualizar status do buffer;
  \item mostrar classe prevista e confianca;
  \item exibir alerta quando Fall for confirmado.
\end{itemize}

\section{Integração IoT: ESP32, Serial e MQTT}

\subsection{ESP32 e Arduino Framework}
O ESP32 executa firmware com:
\begin{itemize}
  \item controle de GPIO para buzzer e LEDs;
  \item recebimento de mensagem de alerta para ativar o alarme local.
\end{itemize}

\subsection{Comunicacao Serial (UART)}
Serial e recomendada para testes:
\begin{itemize}
  \item simples de configurar;
  \item nao exige WiFi;
  \item limita distancia e depende do cabo.
\end{itemize}
No contexto do projeto, usa-se configuracao de porta COM e baudrate.

\subsection{Comunicacao MQTT}
MQTT e um protocolo de mensageria publish/subscribe:
\begin{itemize}
  \item produtores publicam mensagens em um topico;
  \item assinantes recebem mensagens de um topico.
\end{itemize}
No projeto, o topico usado e ``fall\_detection/alerts''.

\subsection{Broker Mosquitto e WebSocket}
O Mosquitto e configurado com:
\begin{itemize}
  \item listener MQTT em 1883;
  \item listener WebSocket em 9001 para permitir conexao do cliente do app.
\end{itemize}

\subsection{Formato de Mensagem em JSON}
O payload do alerta e JSON. Conceitos importantes:
\begin{itemize}
  \item campos de controle: alert, confidence, timestamp;
  \item metadata opcional: frame\_id, model, etc.
\end{itemize}
Essa estrutura facilita interpretacao e depuracao, e permite evolucao do sistema sem quebrar compatibilidade.

\subsection{Estratégias de Robustez}
O estudo deve cobrir:
\begin{itemize}
  \item reconexao ao broker;
  \item tratamento de falhas (timeouts, portas incorretas, firewall);
  \item seguranca (quando usar MQTT com autenticação).
\end{itemize}

\section{App Mobile: React Native + Expo + MQTT over WebSocket}

\subsection{Por que WebSocket para MQTT no App}
Em ambiente JavaScript do React Native, conexoes TCP diretas nem sempre sao praticas. Assim, usa-se MQTT sobre WebSocket (WebSocket listener no broker) para permitir conectividade pelo ambiente mobile.

\subsection{Expo e Expo Go}
Expo simplifica:
\begin{itemize}
  \item configuracao do projeto;
  \item execucao no celular via Expo Go;
  \item acesso a API para audio, vibração e armazenamento.
\end{itemize}

\subsection{Gerenciamento de Estado com Context API}
O App utiliza Context API e useReducer para organizar:
\begin{itemize}
  \item estado de conexao MQTT;
  \item estado do alarme (ativo/inativo);
  \item eventos recebidos e persistidos.
\end{itemize}

\subsection{Servicos (Services)}
Arquitetura separa:
\begin{itemize}
  \item MqttService: conexao e subscribe no topico;
  \item AlarmService: som em loop e vibracao;
  \item EventStorage: persistencia no AsyncStorage.
\end{itemize}

\subsection{Fluxo do Alerta no App}
Fluxo conceitual:
\begin{enumerate}
  \item App recebe mensagem MQTT (JSON).
  \item Valida campo ``alert'' e ``confidence'' vs limiar.
  \item Aciona tela e inicia o alarme (som e vibração).
  \item Permite ao cuidador confirmar ou registrar falso alarme.
  \item Mantem historico persistido e filtro por tipo.
\end{enumerate}

\subsection{Historico e AsyncStorage}
AsyncStorage permite persistir eventos/configs localmente no dispositivo. Conceito central:
\begin{itemize}
  \item persistencia entre sessoes;
  \item consulta e filtros por tipo;
  \item limite de tamanho do historico (ex.: ate 200 eventos).
\end{itemize}

\section{Reprodutibilidade: Logs e Artefatos}

\subsection{Estrutura de Logs do Treinamento}
Cada treino gera uma pasta run-YYYYMMDD-HHMMSS/ contendo:
\begin{itemize}
  \item run\_metadata.json: hiperparametros e configuracoes;
  \item history.json / history.csv: evolucao por epoca;
  \item curves.png: graficos de perda e acuracia;
  \item classification\_report.txt: precision/recall/F1 por classe;
  \item confusion\_matrix.png: matriz visual;
  \item final\_metrics.json: loss e accuracy finais.
\end{itemize}

\subsection{Versionamento de Artefatos}
Boas praticas:
\begin{itemize}
  \item versionar artefatos leves (JSON/CSV/TXT/PNG);
  \item evitar versionar arquivos grandes (.h5 e dados brutos) ou usar Git LFS (se necessario).
\end{itemize}

\section{Limitações e Consideracoes de Ciencia do Mundo Real}

\subsection{Generalizacao}
Resultados excelentes no conjunto de teste podem nao refletir desempenho em:
\begin{itemize}
  \item outros ambientes/cameras;
  \item mudancas de iluminacao;
  \item posicoes e alturas distintas;
  \item variacoes de vestimenta e movimento.
\end{itemize}

\subsection{Falsos Positivos Temporais}
Deteccoes antecipadas podem ocorrer em transicoes de movimento. A mitigacao usa:
\begin{itemize}
  \item limiar de confidence mais alto;
  \item confirmacao por varias predições consecutivas.
\end{itemize}

\subsection{Performance}
Em hardware limitado:
\begin{itemize}
  \item usar TFLite/INT8;
  \item skip frames e buffer reduzido;
  \item evitar render pesado em loops (otimizar visualizacao).
\end{itemize}

\section{Referencias (para voce citar no Overleaf)}
Sugestoes de referencias utilizadas no contexto do projeto:
\begin{thebibliography}{9}
  \bibitem{kwolek2014} Kwolek, B.; Kepski, M. Human fall detection on embedded platform using depth maps and wireless accelerometer. Computer Methods and Programs in Biomedicine, 117(3), 489--501, 2014.
  \bibitem{mobnetv2} Sandler, M. et al. MobileNetV2: Inverted Residuals and Linear Bottlenecks. CVPR, 2018.
  \bibitem{lstm} Hochreiter, S.; Schmidhuber, J. Long Short-Term Memory. Neural Computation, 1997.
  \bibitem{lu2019} Lu, N. et al. Deep Learning for Fall Detection: Three-Dimensional CNN Combined With LSTM on Video Kinematic Data. IEEE JBHI, 2019.
  \bibitem{mubashir2013} Mubashir, M.; Shao, L.; Seed, L. A survey on fall detection: Principles and approaches. Neurocomputing, 2013.
  \bibitem{mqtt} MQTT.org. MQTT v3.1.1 Specification. (Documentacao oficial do protocolo).
  \bibitem{mosquitto} Eclipse Mosquitto. Documentacao e referencia do broker.
\end{thebibliography}


%
% Obs.: texto em formato LaTeX (secoes e subsecoes) para estudo teorico.

\section{Estudo - Tecnologias e Conceitos do Projeto}

\subsection{Visao Geral}
Este projeto implementa um \textbf{sistema de deteccao de quedas} baseado em aprendizado profundo e visao computacional, com integracao de notificacoes via \textbf{ESP32} (hardware) e \textbf{aplicativo mobile} (React Native/Expo). A arquitetura ponta a ponta envolve:
captura de video por camara, processamento em janelas temporais, inferencia com uma rede \textbf{CNN+LSTM}, e distribuicao de alertas usando \textbf{MQTT} (com broker Mosquitto) para acionamento local (buzzer/LEDs) e remoto (celular).

\subsection{Objetivos do Estudo}
Compreender e relacionar:
\begin{itemize}
  \item Tecnicas de \textbf{deep learning} para classificacao de eventos em video.
  \item Processamento de video com \textbf{OpenCV} e preparo de dados (ETL).
  \item Padroes de rede neural: \textbf{MobileNetV2}, \textbf{TimeDistributed}, \textbf{LSTM} e classificacao binaria.
  \item Treinamento supervisionado: hiperparametros, callbacks e regularizacao.
  \item Avaliacao estatistica: matriz de confusao, precision, recall e F1-score.
  \item Inferencia e otimizacao para \textbf{CPU} com \textbf{TFLite} e quantizacao.
  \item Integracao IoT: \textbf{ESP32}, comunicacao \textbf{Serial/UART}, \textbf{MQTT} e mensageria em JSON.
  \item App mobile: \textbf{React Native}, \textbf{Expo}, WebSocket para MQTT e persistencia local com AsyncStorage.
  \item Reprodutibilidade: logs de treinamento e controle de artefatos.
\end{itemize}

\subsection{Stack Tecnologico (Resumo)}
\begin{itemize}
  \item \textbf{Backend IA e Visao Computacional:} Python, TensorFlow/Keras, OpenCV, NumPy, scikit-learn.
  \item \textbf{Treinamento e Avaliacao:} tf.data, callbacks (EarlyStopping, ReduceLROnPlateau, ModelCheckpoint), matplotlib.
  \item \textbf{Inferencia Otimizada:} TFLite com quantizacao INT8 (quando disponivel).
  \item \textbf{IoT:} ESP32 (Arduino), Serial (UART) via pyserial e MQTT via paho-mqtt.
  \item \textbf{Broker:} Eclipse Mosquitto com listener MQTT (1883) e WebSocket (9001) para o app.
  \item \textbf{App Mobile:} React Native + Expo + TypeScript.
  \item \textbf{Persistencia e Media:} AsyncStorage (historico/config), audio em loop e vibra (haptics) via bibliotecas do Expo.
\end{itemize}

\section{Tema Central: Deteccao de Quedas em Video}

\subsection{Por que Video e Tempo Importam}
Quedas sao eventos dinamicos: nao se resumem a uma postura unica. Por isso, e necessario modelar a \textbf{evolucao temporal} do movimento. Em vez de classificar frames isolados, o sistema trabalha com \textbf{sequencias} (janelas temporais) de frames.

\subsection{Abordagem de Classificacao Binaria}
O problema e formulado como classificacao binaria:
\begin{itemize}
  \item \textbf{Normal} (0)
  \item \textbf{Fall} (1)
\end{itemize}
A rede produz um valor no intervalo $[0,1]$ via ativacao \textbf{sigmoid}, interpretado como probabilidade da classe Fall.

\section{Modelos de Deep Learning: CNN+LSTM}

\subsection{Conceito de CNN (Convolutional Neural Network)}
CNNs extraem caracteristicas espaciais relevantes (bordas, texturas e padrões visuais). Em classificacao de imagens, uma CNN mapeia uma entrada de dimensoes $H \times W \times C$ para uma representacao vetorial.

\subsection{Transfer Learning com MobileNetV2}
O projeto utiliza \textbf{MobileNetV2} como extrator de caracteristicas por:
\begin{itemize}
  \item eficiencia computacional (arquiteturas de mobile para menor latencia);
  \item uso de pesos pre-treinados no ImageNet (transfer learning).
\end{itemize}
Transfer learning significa congelar a base inicialmente (ou em parte) para reduzir tempo e necessidade de dados rotulados.

\subsection{TimeDistributed e Processamento de Sequencias}
Para cada sequencia de $T=20$ frames, a CNN e aplicada individualmente a cada frame. A camada \textbf{TimeDistributed} replica a operacao da CNN ao longo do eixo temporal:
\begin{itemize}
  \item entrada: $(T, H, W, C)$ por amostra;
  \item saida: vetores de caracteristicas por frame (para alimentacao temporal).
\end{itemize}

\subsection{LSTM (Long Short-Term Memory)}
A LSTM modela dependencia temporal e aprendizados de como caracteristicas mudam ao longo do tempo. O componente LSTM tende a lidar com informacoes de longo prazo atraves de \textbf{portas de controle} (input/forget/output) e estados internos.

\subsection{Regularizacao (Dropout)}
O modelo aplica \textbf{Dropout} na camada de saida da LSTM para reduzir overfitting. Em geral, Dropout desativa aleatoriamente parte das unidades durante o treinamento, forçando o modelo a aprender representacoes mais robustas.

\subsection{Funcao de Perda e Otimizador}
O projeto utiliza:
\begin{itemize}
  \item \textbf{Loss:} binary crossentropy (adequada a classificacao binaria);
  \item \textbf{Otimizador:} Adam com learning rate configurado.
\end{itemize}

\section{Preparacao de Dados (ETL) para Video}

\subsection{Dataset UR Fall Detection}
O dataset usado contem:
\begin{itemize}
  \item sequencias de quedas simuladas (Fall);
  \item atividades diarias (ADL/Normal).
\end{itemize}
O sistema assume que videos podem ser organizados por pastas/classe.

\subsection{Pipeline de Pre-processamento}
O fluxo de preparo de dados envolve:
\begin{itemize}
  \item identificar arquivos por busca recursiva;
  \item extrair frames ou converter sequencias para um formato de video padronizado;
  \item redimensionar imagens para o tamanho exigido pelo modelo;
  \item normalizar valores (ex.: escala para $[0,1]$);
  \item construir amostras via \textbf{sliding window} (janelas temporais sobre o video).
\end{itemize}

\subsection{Sliding Window e Passo Temporal}
O sliding window cria varias amostras a partir de um unico video. O passo temporal (stride) controla a sobreposicao entre janelas:
\begin{itemize}
  \item stride menor => mais amostras, maior custo computacional;
  \item stride maior => menos amostras, potencialmente menos cobertura temporal.
\end{itemize}

\subsection{Data Augmentation}
Para melhorar generalizacao, o projeto aplica aumentos aleatorios em tempo real em sequencias, como:
\begin{itemize}
  \item flip horizontal;
  \item variacao de brilho.
\end{itemize}
Essa etapa reduz overfitting e melhora robustez para variacoes do ambiente.

\subsection{Split Estratificado (Treino/Teste)}
A separacao treino/teste e estratificada para preservar proporcao de classes em cada subconjunto. Isso evita que o modelo seja treinado com uma distribuicao distorcida.

\section{Treinamento do Modelo}

\subsection{Hiperparametros}
Aspectos estudaveis incluem:
\begin{itemize}
  \item numero de epocas;
  \item batch size;
  \item learning rate;
  \item tamanho da sequencia (T=20);
  \item dimensao de entrada (H=W=128).
\end{itemize}

\subsection{tf.data e Eficiência}
O uso de tf.data ajuda a:
\begin{itemize}
  \item prefetch e pipelining de dados;
  \item reduz overhead de carregamento;
  \item suportar augmentation em tempo real.
\end{itemize}

\subsection{Callbacks}
Callbacks implementados:
\begin{itemize}
  \item \textbf{ModelCheckpoint:} salva o melhor modelo conforme val\_accuracy;
  \item \textbf{EarlyStopping:} interrompe treinamento quando val\_loss para de melhorar;
  \item \textbf{ReduceLROnPlateau:} reduz learning rate quando val\_loss estagna.
\end{itemize}

\subsection{Mixed Precision (Opcional)}
Quando habilitado, mixed precision acelera treinamento em hardware compatível, reduzindo memoria e aumentando throughput. No projeto, o modo esta configurado como desativado por padrao.

\section{Avaliacao: Metrics e Analise de Erros}

\subsection{Accuracy e Loss}
Accuracy mede a proporcao de amostras corretamente classificadas. Loss quantifica o erro durante treinamento e avaliacao, sendo util para observar convergencia.

\subsection{Precision, Recall e F1-score}
Precision (para a classe prevista) e Recall (para a classe real) sao essenciais em cenarios desbalanceados. F1-score combina ambos em uma unica medida.

\subsection{Matriz de Confusao}
A matriz de confusao $C$ resume:
\begin{itemize}
  \item True Positives (acertos de Fall);
  \item True Negatives (acertos de Normal);
  \item False Positives (Normal previsto como Fall);
  \item False Negatives (Fall previsto como Normal).
\end{itemize}
Ela e a principal fonte visual de entendimento de erros.

\subsection{Leitura dos Graficos}
Os graficos de curves.png mostram:
\begin{itemize}
  \item train loss vs val loss;
  \item train accuracy vs val accuracy.
\end{itemize}
Quando val e train evoluem de forma semelhante, tende a indicar boa generalizacao (sem overfitting significativo).

\section{Inferencia em Tempo Real e Otimizacao}

\subsection{Buffer Circular de Sequencias}
Durante inferencia, o sistema acumula frames em um buffer ate completar T=20. A cada passo de inferencia, a janela atual e processada.

\subsection{Skip Frames}
O uso de inferencia a cada N frames (skip frames) reduz custo computacional e viabiliza execucao em CPU.

\subsection{Limiar e Confirmacao Temporal}
Para reduzir falsos positivos temporais, aplica-se:
\begin{itemize}
  \item um limiar mais conservador para considerar Fall;
  \item confirmacao temporal: so dispara Fall apos N predições consecutivas.
\end{itemize}

\subsection{TFLite e Quantizacao INT8}
O modelo e convertido para TFLite para desempenho. Quantizacao INT8 reduz tamanho e acelera em CPU, com possivel impacto na precisao (avaliar sempre).

\section{Visao Computacional com OpenCV}

\subsection{Captura de Video}
OpenCV utiliza VideoCapture para:
\begin{itemize}
  \item acessar webcam (indice 0);
  \item acessar arquivo de video (caminho do arquivo).
\end{itemize}

\subsection{Redimensionamento e Normalizacao}
A entrada do modelo precisa ser padronizada:
\begin{itemize}
  \item resize para HxW;
  \item conversao para float32 e normalizacao para [0,1];
  \item empilhamento temporal em formato correto para a rede.
\end{itemize}

\subsection{Visualizacao}
O sistema usa overlays com texto e paines para:
\begin{itemize}
  \item visualizar status do buffer;
  \item mostrar classe prevista e confianca;
  \item exibir alerta quando Fall for confirmado.
\end{itemize}

\section{Integração IoT: ESP32, Serial e MQTT}

\subsection{ESP32 e Arduino Framework}
O ESP32 executa firmware com:
\begin{itemize}
  \item controle de GPIO para buzzer e LEDs;
  \item recebimento de mensagem de alerta para ativar o alarme local.
\end{itemize}

\subsection{Comunicacao Serial (UART)}
Serial e recomendada para testes:
\begin{itemize}
  \item simples de configurar;
  \item nao exige WiFi;
  \item limita distancia e depende do cabo.
\end{itemize}
No contexto do projeto, usa-se configuracao de porta COM e baudrate.

\subsection{Comunicacao MQTT}
MQTT e um protocolo de mensageria publish/subscribe:
\begin{itemize}
  \item produtores publicam mensagens em um topico;
  \item assinantes recebem mensagens de um topico.
\end{itemize}
No projeto, o topico usado e ``fall\_detection/alerts''.

\subsection{Broker Mosquitto e WebSocket}
O Mosquitto e configurado com:
\begin{itemize}
  \item listener MQTT em 1883;
  \item listener WebSocket em 9001 para permitir conexao do cliente do app.
\end{itemize}

\subsection{Formato de Mensagem em JSON}
O payload do alerta e JSON. Conceitos importantes:
\begin{itemize}
  \item campos de controle: alert, confidence, timestamp;
  \item metadata opcional: frame\_id, model, etc.
\end{itemize}
Essa estrutura facilita interpretacao e depuracao, e permite evolucao do sistema sem quebrar compatibilidade.

\subsection{Estratégias de Robustez}
O estudo deve cobrir:
\begin{itemize}
  \item reconexao ao broker;
  \item tratamento de falhas (timeouts, portas incorretas, firewall);
  \item seguranca (quando usar MQTT com autenticação).
\end{itemize}

\section{App Mobile: React Native + Expo + MQTT over WebSocket}

\subsection{Por que WebSocket para MQTT no App}
Em ambiente JavaScript do React Native, conexoes TCP diretas nem sempre sao praticas. Assim, usa-se MQTT sobre WebSocket (WebSocket listener no broker) para permitir conectividade pelo ambiente mobile.

\subsection{Expo e Expo Go}
Expo simplifica:
\begin{itemize}
  \item configuracao do projeto;
  \item execucao no celular via Expo Go;
  \item acesso a API para audio, vibração e armazenamento.
\end{itemize}

\subsection{Gerenciamento de Estado com Context API}
O App utiliza Context API e useReducer para organizar:
\begin{itemize}
  \item estado de conexao MQTT;
  \item estado do alarme (ativo/inativo);
  \item eventos recebidos e persistidos.
\end{itemize}

\subsection{Servicos (Services)}
Arquitetura separa:
\begin{itemize}
  \item MqttService: conexao e subscribe no topico;
  \item AlarmService: som em loop e vibracao;
  \item EventStorage: persistencia no AsyncStorage.
\end{itemize}

\subsection{Fluxo do Alerta no App}
Fluxo conceitual:
\begin{enumerate}
  \item App recebe mensagem MQTT (JSON).
  \item Valida campo ``alert'' e ``confidence'' vs limiar.
  \item Aciona tela e inicia o alarme (som e vibração).
  \item Permite ao cuidador confirmar ou registrar falso alarme.
  \item Mantem historico persistido e filtro por tipo.
\end{enumerate}

\subsection{Historico e AsyncStorage}
AsyncStorage permite persistir eventos/configs localmente no dispositivo. Conceito central:
\begin{itemize}
  \item persistencia entre sessoes;
  \item consulta e filtros por tipo;
  \item limite de tamanho do historico (ex.: ate 200 eventos).
\end{itemize}

\section{Reprodutibilidade: Logs e Artefatos}

\subsection{Estrutura de Logs do Treinamento}
Cada treino gera uma pasta run-YYYYMMDD-HHMMSS/ contendo:
\begin{itemize}
  \item run\_metadata.json: hiperparametros e configuracoes;
  \item history.json / history.csv: evolucao por epoca;
  \item curves.png: graficos de perda e acuracia;
  \item classification\_report.txt: precision/recall/F1 por classe;
  \item confusion\_matrix.png: matriz visual;
  \item final\_metrics.json: loss e accuracy finais.
\end{itemize}

\subsection{Versionamento de Artefatos}
Boas praticas:
\begin{itemize}
  \item versionar artefatos leves (JSON/CSV/TXT/PNG);
  \item evitar versionar arquivos grandes (.h5 e dados brutos) ou usar Git LFS (se necessario).
\end{itemize}

\section{Limitações e Consideracoes de Ciencia do Mundo Real}

\subsection{Generalizacao}
Resultados excelentes no conjunto de teste podem nao refletir desempenho em:
\begin{itemize}
  \item outros ambientes/cameras;
  \item mudancas de iluminacao;
  \item posicoes e alturas distintas;
  \item variacoes de vestimenta e movimento.
\end{itemize}

\subsection{Falsos Positivos Temporais}
Deteccoes antecipadas podem ocorrer em transicoes de movimento. A mitigacao usa:
\begin{itemize}
  \item limiar de confidence mais alto;
  \item confirmacao por varias predições consecutivas.
\end{itemize}

\subsection{Performance}
Em hardware limitado:
\begin{itemize}
  \item usar TFLite/INT8;
  \item skip frames e buffer reduzido;
  \item evitar render pesado em loops (otimizar visualizacao).
\end{itemize}

\section{Referencias (para voce citar no Overleaf)}
Sugestoes de referencias utilizadas no contexto do projeto:
\begin{thebibliography}{9}
  \bibitem{kwolek2014} Kwolek, B.; Kepski, M. Human fall detection on embedded platform using depth maps and wireless accelerometer. Computer Methods and Programs in Biomedicine, 117(3), 489--501, 2014.
  \bibitem{mobnetv2} Sandler, M. et al. MobileNetV2: Inverted Residuals and Linear Bottlenecks. CVPR, 2018.
  \bibitem{lstm} Hochreiter, S.; Schmidhuber, J. Long Short-Term Memory. Neural Computation, 1997.
  \bibitem{lu2019} Lu, N. et al. Deep Learning for Fall Detection: Three-Dimensional CNN Combined With LSTM on Video Kinematic Data. IEEE JBHI, 2019.
  \bibitem{mubashir2013} Mubashir, M.; Shao, L.; Seed, L. A survey on fall detection: Principles and approaches. Neurocomputing, 2013.
  \bibitem{mqtt} MQTT.org. MQTT v3.1.1 Specification. (Documentacao oficial do protocolo).
  \bibitem{mosquitto} Eclipse Mosquitto. Documentacao e referencia do broker.
\end{thebibliography}


%
% Obs.: texto em formato LaTeX (secoes e subsecoes) para estudo teorico.

\section{Estudo - Tecnologias e Conceitos do Projeto}

\subsection{Visao Geral}
Este projeto implementa um \textbf{sistema de deteccao de quedas} baseado em aprendizado profundo e visao computacional, com integracao de notificacoes via \textbf{ESP32} (hardware) e \textbf{aplicativo mobile} (React Native/Expo). A arquitetura ponta a ponta envolve:
captura de video por camara, processamento em janelas temporais, inferencia com uma rede \textbf{CNN+LSTM}, e distribuicao de alertas usando \textbf{MQTT} (com broker Mosquitto) para acionamento local (buzzer/LEDs) e remoto (celular).

\subsection{Objetivos do Estudo}
Compreender e relacionar:
\begin{itemize}
  \item Tecnicas de \textbf{deep learning} para classificacao de eventos em video.
  \item Processamento de video com \textbf{OpenCV} e preparo de dados (ETL).
  \item Padroes de rede neural: \textbf{MobileNetV2}, \textbf{TimeDistributed}, \textbf{LSTM} e classificacao binaria.
  \item Treinamento supervisionado: hiperparametros, callbacks e regularizacao.
  \item Avaliacao estatistica: matriz de confusao, precision, recall e F1-score.
  \item Inferencia e otimizacao para \textbf{CPU} com \textbf{TFLite} e quantizacao.
  \item Integracao IoT: \textbf{ESP32}, comunicacao \textbf{Serial/UART}, \textbf{MQTT} e mensageria em JSON.
  \item App mobile: \textbf{React Native}, \textbf{Expo}, WebSocket para MQTT e persistencia local com AsyncStorage.
  \item Reprodutibilidade: logs de treinamento e controle de artefatos.
\end{itemize}

\subsection{Stack Tecnologico (Resumo)}
\begin{itemize}
  \item \textbf{Backend IA e Visao Computacional:} Python, TensorFlow/Keras, OpenCV, NumPy, scikit-learn.
  \item \textbf{Treinamento e Avaliacao:} tf.data, callbacks (EarlyStopping, ReduceLROnPlateau, ModelCheckpoint), matplotlib.
  \item \textbf{Inferencia Otimizada:} TFLite com quantizacao INT8 (quando disponivel).
  \item \textbf{IoT:} ESP32 (Arduino), Serial (UART) via pyserial e MQTT via paho-mqtt.
  \item \textbf{Broker:} Eclipse Mosquitto com listener MQTT (1883) e WebSocket (9001) para o app.
  \item \textbf{App Mobile:} React Native + Expo + TypeScript.
  \item \textbf{Persistencia e Media:} AsyncStorage (historico/config), audio em loop e vibra (haptics) via bibliotecas do Expo.
\end{itemize}

\section{Tema Central: Deteccao de Quedas em Video}

\subsection{Por que Video e Tempo Importam}
Quedas sao eventos dinamicos: nao se resumem a uma postura unica. Por isso, e necessario modelar a \textbf{evolucao temporal} do movimento. Em vez de classificar frames isolados, o sistema trabalha com \textbf{sequencias} (janelas temporais) de frames.

\subsection{Abordagem de Classificacao Binaria}
O problema e formulado como classificacao binaria:
\begin{itemize}
  \item \textbf{Normal} (0)
  \item \textbf{Fall} (1)
\end{itemize}
A rede produz um valor no intervalo $[0,1]$ via ativacao \textbf{sigmoid}, interpretado como probabilidade da classe Fall.

\section{Modelos de Deep Learning: CNN+LSTM}

\subsection{Conceito de CNN (Convolutional Neural Network)}
CNNs extraem caracteristicas espaciais relevantes (bordas, texturas e padrões visuais). Em classificacao de imagens, uma CNN mapeia uma entrada de dimensoes $H \times W \times C$ para uma representacao vetorial.

\subsection{Transfer Learning com MobileNetV2}
O projeto utiliza \textbf{MobileNetV2} como extrator de caracteristicas por:
\begin{itemize}
  \item eficiencia computacional (arquiteturas de mobile para menor latencia);
  \item uso de pesos pre-treinados no ImageNet (transfer learning).
\end{itemize}
Transfer learning significa congelar a base inicialmente (ou em parte) para reduzir tempo e necessidade de dados rotulados.

\subsection{TimeDistributed e Processamento de Sequencias}
Para cada sequencia de $T=20$ frames, a CNN e aplicada individualmente a cada frame. A camada \textbf{TimeDistributed} replica a operacao da CNN ao longo do eixo temporal:
\begin{itemize}
  \item entrada: $(T, H, W, C)$ por amostra;
  \item saida: vetores de caracteristicas por frame (para alimentacao temporal).
\end{itemize}

\subsection{LSTM (Long Short-Term Memory)}
A LSTM modela dependencia temporal e aprendizados de como caracteristicas mudam ao longo do tempo. O componente LSTM tende a lidar com informacoes de longo prazo atraves de \textbf{portas de controle} (input/forget/output) e estados internos.

\subsection{Regularizacao (Dropout)}
O modelo aplica \textbf{Dropout} na camada de saida da LSTM para reduzir overfitting. Em geral, Dropout desativa aleatoriamente parte das unidades durante o treinamento, forçando o modelo a aprender representacoes mais robustas.

\subsection{Funcao de Perda e Otimizador}
O projeto utiliza:
\begin{itemize}
  \item \textbf{Loss:} binary crossentropy (adequada a classificacao binaria);
  \item \textbf{Otimizador:} Adam com learning rate configurado.
\end{itemize}

\section{Preparacao de Dados (ETL) para Video}

\subsection{Dataset UR Fall Detection}
O dataset usado contem:
\begin{itemize}
  \item sequencias de quedas simuladas (Fall);
  \item atividades diarias (ADL/Normal).
\end{itemize}
O sistema assume que videos podem ser organizados por pastas/classe.

\subsection{Pipeline de Pre-processamento}
O fluxo de preparo de dados envolve:
\begin{itemize}
  \item identificar arquivos por busca recursiva;
  \item extrair frames ou converter sequencias para um formato de video padronizado;
  \item redimensionar imagens para o tamanho exigido pelo modelo;
  \item normalizar valores (ex.: escala para $[0,1]$);
  \item construir amostras via \textbf{sliding window} (janelas temporais sobre o video).
\end{itemize}

\subsection{Sliding Window e Passo Temporal}
O sliding window cria varias amostras a partir de um unico video. O passo temporal (stride) controla a sobreposicao entre janelas:
\begin{itemize}
  \item stride menor => mais amostras, maior custo computacional;
  \item stride maior => menos amostras, potencialmente menos cobertura temporal.
\end{itemize}

\subsection{Data Augmentation}
Para melhorar generalizacao, o projeto aplica aumentos aleatorios em tempo real em sequencias, como:
\begin{itemize}
  \item flip horizontal;
  \item variacao de brilho.
\end{itemize}
Essa etapa reduz overfitting e melhora robustez para variacoes do ambiente.

\subsection{Split Estratificado (Treino/Teste)}
A separacao treino/teste e estratificada para preservar proporcao de classes em cada subconjunto. Isso evita que o modelo seja treinado com uma distribuicao distorcida.

\section{Treinamento do Modelo}

\subsection{Hiperparametros}
Aspectos estudaveis incluem:
\begin{itemize}
  \item numero de epocas;
  \item batch size;
  \item learning rate;
  \item tamanho da sequencia (T=20);
  \item dimensao de entrada (H=W=128).
\end{itemize}

\subsection{tf.data e Eficiência}
O uso de tf.data ajuda a:
\begin{itemize}
  \item prefetch e pipelining de dados;
  \item reduz overhead de carregamento;
  \item suportar augmentation em tempo real.
\end{itemize}

\subsection{Callbacks}
Callbacks implementados:
\begin{itemize}
  \item \textbf{ModelCheckpoint:} salva o melhor modelo conforme val\_accuracy;
  \item \textbf{EarlyStopping:} interrompe treinamento quando val\_loss para de melhorar;
  \item \textbf{ReduceLROnPlateau:} reduz learning rate quando val\_loss estagna.
\end{itemize}

\subsection{Mixed Precision (Opcional)}
Quando habilitado, mixed precision acelera treinamento em hardware compatível, reduzindo memoria e aumentando throughput. No projeto, o modo esta configurado como desativado por padrao.

\section{Avaliacao: Metrics e Analise de Erros}

\subsection{Accuracy e Loss}
Accuracy mede a proporcao de amostras corretamente classificadas. Loss quantifica o erro durante treinamento e avaliacao, sendo util para observar convergencia.

\subsection{Precision, Recall e F1-score}
Precision (para a classe prevista) e Recall (para a classe real) sao essenciais em cenarios desbalanceados. F1-score combina ambos em uma unica medida.

\subsection{Matriz de Confusao}
A matriz de confusao $C$ resume:
\begin{itemize}
  \item True Positives (acertos de Fall);
  \item True Negatives (acertos de Normal);
  \item False Positives (Normal previsto como Fall);
  \item False Negatives (Fall previsto como Normal).
\end{itemize}
Ela e a principal fonte visual de entendimento de erros.

\subsection{Leitura dos Graficos}
Os graficos de curves.png mostram:
\begin{itemize}
  \item train loss vs val loss;
  \item train accuracy vs val accuracy.
\end{itemize}
Quando val e train evoluem de forma semelhante, tende a indicar boa generalizacao (sem overfitting significativo).

\section{Inferencia em Tempo Real e Otimizacao}

\subsection{Buffer Circular de Sequencias}
Durante inferencia, o sistema acumula frames em um buffer ate completar T=20. A cada passo de inferencia, a janela atual e processada.

\subsection{Skip Frames}
O uso de inferencia a cada N frames (skip frames) reduz custo computacional e viabiliza execucao em CPU.

\subsection{Limiar e Confirmacao Temporal}
Para reduzir falsos positivos temporais, aplica-se:
\begin{itemize}
  \item um limiar mais conservador para considerar Fall;
  \item confirmacao temporal: so dispara Fall apos N predições consecutivas.
\end{itemize}

\subsection{TFLite e Quantizacao INT8}
O modelo e convertido para TFLite para desempenho. Quantizacao INT8 reduz tamanho e acelera em CPU, com possivel impacto na precisao (avaliar sempre).

\section{Visao Computacional com OpenCV}

\subsection{Captura de Video}
OpenCV utiliza VideoCapture para:
\begin{itemize}
  \item acessar webcam (indice 0);
  \item acessar arquivo de video (caminho do arquivo).
\end{itemize}

\subsection{Redimensionamento e Normalizacao}
A entrada do modelo precisa ser padronizada:
\begin{itemize}
  \item resize para HxW;
  \item conversao para float32 e normalizacao para [0,1];
  \item empilhamento temporal em formato correto para a rede.
\end{itemize}

\subsection{Visualizacao}
O sistema usa overlays com texto e paines para:
\begin{itemize}
  \item visualizar status do buffer;
  \item mostrar classe prevista e confianca;
  \item exibir alerta quando Fall for confirmado.
\end{itemize}

\section{Integração IoT: ESP32, Serial e MQTT}

\subsection{ESP32 e Arduino Framework}
O ESP32 executa firmware com:
\begin{itemize}
  \item controle de GPIO para buzzer e LEDs;
  \item recebimento de mensagem de alerta para ativar o alarme local.
\end{itemize}

\subsection{Comunicacao Serial (UART)}
Serial e recomendada para testes:
\begin{itemize}
  \item simples de configurar;
  \item nao exige WiFi;
  \item limita distancia e depende do cabo.
\end{itemize}
No contexto do projeto, usa-se configuracao de porta COM e baudrate.

\subsection{Comunicacao MQTT}
MQTT e um protocolo de mensageria publish/subscribe:
\begin{itemize}
  \item produtores publicam mensagens em um topico;
  \item assinantes recebem mensagens de um topico.
\end{itemize}
No projeto, o topico usado e ``fall\_detection/alerts''.

\subsection{Broker Mosquitto e WebSocket}
O Mosquitto e configurado com:
\begin{itemize}
  \item listener MQTT em 1883;
  \item listener WebSocket em 9001 para permitir conexao do cliente do app.
\end{itemize}

\subsection{Formato de Mensagem em JSON}
O payload do alerta e JSON. Conceitos importantes:
\begin{itemize}
  \item campos de controle: alert, confidence, timestamp;
  \item metadata opcional: frame\_id, model, etc.
\end{itemize}
Essa estrutura facilita interpretacao e depuracao, e permite evolucao do sistema sem quebrar compatibilidade.

\subsection{Estratégias de Robustez}
O estudo deve cobrir:
\begin{itemize}
  \item reconexao ao broker;
  \item tratamento de falhas (timeouts, portas incorretas, firewall);
  \item seguranca (quando usar MQTT com autenticação).
\end{itemize}

\section{App Mobile: React Native + Expo + MQTT over WebSocket}

\subsection{Por que WebSocket para MQTT no App}
Em ambiente JavaScript do React Native, conexoes TCP diretas nem sempre sao praticas. Assim, usa-se MQTT sobre WebSocket (WebSocket listener no broker) para permitir conectividade pelo ambiente mobile.

\subsection{Expo e Expo Go}
Expo simplifica:
\begin{itemize}
  \item configuracao do projeto;
  \item execucao no celular via Expo Go;
  \item acesso a API para audio, vibração e armazenamento.
\end{itemize}

\subsection{Gerenciamento de Estado com Context API}
O App utiliza Context API e useReducer para organizar:
\begin{itemize}
  \item estado de conexao MQTT;
  \item estado do alarme (ativo/inativo);
  \item eventos recebidos e persistidos.
\end{itemize}

\subsection{Servicos (Services)}
Arquitetura separa:
\begin{itemize}
  \item MqttService: conexao e subscribe no topico;
  \item AlarmService: som em loop e vibracao;
  \item EventStorage: persistencia no AsyncStorage.
\end{itemize}

\subsection{Fluxo do Alerta no App}
Fluxo conceitual:
\begin{enumerate}
  \item App recebe mensagem MQTT (JSON).
  \item Valida campo ``alert'' e ``confidence'' vs limiar.
  \item Aciona tela e inicia o alarme (som e vibração).
  \item Permite ao cuidador confirmar ou registrar falso alarme.
  \item Mantem historico persistido e filtro por tipo.
\end{enumerate}

\subsection{Historico e AsyncStorage}
AsyncStorage permite persistir eventos/configs localmente no dispositivo. Conceito central:
\begin{itemize}
  \item persistencia entre sessoes;
  \item consulta e filtros por tipo;
  \item limite de tamanho do historico (ex.: ate 200 eventos).
\end{itemize}

\section{Reprodutibilidade: Logs e Artefatos}

\subsection{Estrutura de Logs do Treinamento}
Cada treino gera uma pasta run-YYYYMMDD-HHMMSS/ contendo:
\begin{itemize}
  \item run\_metadata.json: hiperparametros e configuracoes;
  \item history.json / history.csv: evolucao por epoca;
  \item curves.png: graficos de perda e acuracia;
  \item classification\_report.txt: precision/recall/F1 por classe;
  \item confusion\_matrix.png: matriz visual;
  \item final\_metrics.json: loss e accuracy finais.
\end{itemize}

\subsection{Versionamento de Artefatos}
Boas praticas:
\begin{itemize}
  \item versionar artefatos leves (JSON/CSV/TXT/PNG);
  \item evitar versionar arquivos grandes (.h5 e dados brutos) ou usar Git LFS (se necessario).
\end{itemize}

\section{Limitações e Consideracoes de Ciencia do Mundo Real}

\subsection{Generalizacao}
Resultados excelentes no conjunto de teste podem nao refletir desempenho em:
\begin{itemize}
  \item outros ambientes/cameras;
  \item mudancas de iluminacao;
  \item posicoes e alturas distintas;
  \item variacoes de vestimenta e movimento.
\end{itemize}

\subsection{Falsos Positivos Temporais}
Deteccoes antecipadas podem ocorrer em transicoes de movimento. A mitigacao usa:
\begin{itemize}
  \item limiar de confidence mais alto;
  \item confirmacao por varias predições consecutivas.
\end{itemize}

\subsection{Performance}
Em hardware limitado:
\begin{itemize}
  \item usar TFLite/INT8;
  \item skip frames e buffer reduzido;
  \item evitar render pesado em loops (otimizar visualizacao).
\end{itemize}

\section{Referencias (para voce citar no Overleaf)}
Sugestoes de referencias utilizadas no contexto do projeto:
\begin{thebibliography}{9}
  \bibitem{kwolek2014} Kwolek, B.; Kepski, M. Human fall detection on embedded platform using depth maps and wireless accelerometer. Computer Methods and Programs in Biomedicine, 117(3), 489--501, 2014.
  \bibitem{mobnetv2} Sandler, M. et al. MobileNetV2: Inverted Residuals and Linear Bottlenecks. CVPR, 2018.
  \bibitem{lstm} Hochreiter, S.; Schmidhuber, J. Long Short-Term Memory. Neural Computation, 1997.
  \bibitem{lu2019} Lu, N. et al. Deep Learning for Fall Detection: Three-Dimensional CNN Combined With LSTM on Video Kinematic Data. IEEE JBHI, 2019.
  \bibitem{mubashir2013} Mubashir, M.; Shao, L.; Seed, L. A survey on fall detection: Principles and approaches. Neurocomputing, 2013.
  \bibitem{mqtt} MQTT.org. MQTT v3.1.1 Specification. (Documentacao oficial do protocolo).
  \bibitem{mosquitto} Eclipse Mosquitto. Documentacao e referencia do broker.
\end{thebibliography}


%
% Obs.: texto em formato LaTeX (secoes e subsecoes) para estudo teorico.

\section{Estudo - Tecnologias e Conceitos do Projeto}

\subsection{Visao Geral}
Este projeto implementa um \textbf{sistema de deteccao de quedas} baseado em aprendizado profundo e visao computacional, com integracao de notificacoes via \textbf{ESP32} (hardware) e \textbf{aplicativo mobile} (React Native/Expo). A arquitetura ponta a ponta envolve:
captura de video por camara, processamento em janelas temporais, inferencia com uma rede \textbf{CNN+LSTM}, e distribuicao de alertas usando \textbf{MQTT} (com broker Mosquitto) para acionamento local (buzzer/LEDs) e remoto (celular).

\subsection{Objetivos do Estudo}
Compreender e relacionar:
\begin{itemize}
  \item Tecnicas de \textbf{deep learning} para classificacao de eventos em video.
  \item Processamento de video com \textbf{OpenCV} e preparo de dados (ETL).
  \item Padroes de rede neural: \textbf{MobileNetV2}, \textbf{TimeDistributed}, \textbf{LSTM} e classificacao binaria.
  \item Treinamento supervisionado: hiperparametros, callbacks e regularizacao.
  \item Avaliacao estatistica: matriz de confusao, precision, recall e F1-score.
  \item Inferencia e otimizacao para \textbf{CPU} com \textbf{TFLite} e quantizacao.
  \item Integracao IoT: \textbf{ESP32}, comunicacao \textbf{Serial/UART}, \textbf{MQTT} e mensageria em JSON.
  \item App mobile: \textbf{React Native}, \textbf{Expo}, WebSocket para MQTT e persistencia local com AsyncStorage.
  \item Reprodutibilidade: logs de treinamento e controle de artefatos.
\end{itemize}

\subsection{Stack Tecnologico (Resumo)}
\begin{itemize}
  \item \textbf{Backend IA e Visao Computacional:} Python, TensorFlow/Keras, OpenCV, NumPy, scikit-learn.
  \item \textbf{Treinamento e Avaliacao:} tf.data, callbacks (EarlyStopping, ReduceLROnPlateau, ModelCheckpoint), matplotlib.
  \item \textbf{Inferencia Otimizada:} TFLite com quantizacao INT8 (quando disponivel).
  \item \textbf{IoT:} ESP32 (Arduino), Serial (UART) via pyserial e MQTT via paho-mqtt.
  \item \textbf{Broker:} Eclipse Mosquitto com listener MQTT (1883) e WebSocket (9001) para o app.
  \item \textbf{App Mobile:} React Native + Expo + TypeScript.
  \item \textbf{Persistencia e Media:} AsyncStorage (historico/config), audio em loop e vibra (haptics) via bibliotecas do Expo.
\end{itemize}

\section{Tema Central: Deteccao de Quedas em Video}

\subsection{Por que Video e Tempo Importam}
Quedas sao eventos dinamicos: nao se resumem a uma postura unica. Por isso, e necessario modelar a \textbf{evolucao temporal} do movimento. Em vez de classificar frames isolados, o sistema trabalha com \textbf{sequencias} (janelas temporais) de frames.

\subsection{Abordagem de Classificacao Binaria}
O problema e formulado como classificacao binaria:
\begin{itemize}
  \item \textbf{Normal} (0)
  \item \textbf{Fall} (1)
\end{itemize}
A rede produz um valor no intervalo $[0,1]$ via ativacao \textbf{sigmoid}, interpretado como probabilidade da classe Fall.

\section{Modelos de Deep Learning: CNN+LSTM}

\subsection{Conceito de CNN (Convolutional Neural Network)}
CNNs extraem caracteristicas espaciais relevantes (bordas, texturas e padrões visuais). Em classificacao de imagens, uma CNN mapeia uma entrada de dimensoes $H \times W \times C$ para uma representacao vetorial.

\subsection{Transfer Learning com MobileNetV2}
O projeto utiliza \textbf{MobileNetV2} como extrator de caracteristicas por:
\begin{itemize}
  \item eficiencia computacional (arquiteturas de mobile para menor latencia);
  \item uso de pesos pre-treinados no ImageNet (transfer learning).
\end{itemize}
Transfer learning significa congelar a base inicialmente (ou em parte) para reduzir tempo e necessidade de dados rotulados.

\subsection{TimeDistributed e Processamento de Sequencias}
Para cada sequencia de $T=20$ frames, a CNN e aplicada individualmente a cada frame. A camada \textbf{TimeDistributed} replica a operacao da CNN ao longo do eixo temporal:
\begin{itemize}
  \item entrada: $(T, H, W, C)$ por amostra;
  \item saida: vetores de caracteristicas por frame (para alimentacao temporal).
\end{itemize}

\subsection{LSTM (Long Short-Term Memory)}
A LSTM modela dependencia temporal e aprendizados de como caracteristicas mudam ao longo do tempo. O componente LSTM tende a lidar com informacoes de longo prazo atraves de \textbf{portas de controle} (input/forget/output) e estados internos.

\subsection{Regularizacao (Dropout)}
O modelo aplica \textbf{Dropout} na camada de saida da LSTM para reduzir overfitting. Em geral, Dropout desativa aleatoriamente parte das unidades durante o treinamento, forçando o modelo a aprender representacoes mais robustas.

\subsection{Funcao de Perda e Otimizador}
O projeto utiliza:
\begin{itemize}
  \item \textbf{Loss:} binary crossentropy (adequada a classificacao binaria);
  \item \textbf{Otimizador:} Adam com learning rate configurado.
\end{itemize}

\section{Preparacao de Dados (ETL) para Video}

\subsection{Dataset UR Fall Detection}
O dataset usado contem:
\begin{itemize}
  \item sequencias de quedas simuladas (Fall);
  \item atividades diarias (ADL/Normal).
\end{itemize}
O sistema assume que videos podem ser organizados por pastas/classe.

\subsection{Pipeline de Pre-processamento}
O fluxo de preparo de dados envolve:
\begin{itemize}
  \item identificar arquivos por busca recursiva;
  \item extrair frames ou converter sequencias para um formato de video padronizado;
  \item redimensionar imagens para o tamanho exigido pelo modelo;
  \item normalizar valores (ex.: escala para $[0,1]$);
  \item construir amostras via \textbf{sliding window} (janelas temporais sobre o video).
\end{itemize}

\subsection{Sliding Window e Passo Temporal}
O sliding window cria varias amostras a partir de um unico video. O passo temporal (stride) controla a sobreposicao entre janelas:
\begin{itemize}
  \item stride menor => mais amostras, maior custo computacional;
  \item stride maior => menos amostras, potencialmente menos cobertura temporal.
\end{itemize}

\subsection{Data Augmentation}
Para melhorar generalizacao, o projeto aplica aumentos aleatorios em tempo real em sequencias, como:
\begin{itemize}
  \item flip horizontal;
  \item variacao de brilho.
\end{itemize}
Essa etapa reduz overfitting e melhora robustez para variacoes do ambiente.

\subsection{Split Estratificado (Treino/Teste)}
A separacao treino/teste e estratificada para preservar proporcao de classes em cada subconjunto. Isso evita que o modelo seja treinado com uma distribuicao distorcida.

\section{Treinamento do Modelo}

\subsection{Hiperparametros}
Aspectos estudaveis incluem:
\begin{itemize}
  \item numero de epocas;
  \item batch size;
  \item learning rate;
  \item tamanho da sequencia (T=20);
  \item dimensao de entrada (H=W=128).
\end{itemize}

\subsection{tf.data e Eficiência}
O uso de tf.data ajuda a:
\begin{itemize}
  \item prefetch e pipelining de dados;
  \item reduz overhead de carregamento;
  \item suportar augmentation em tempo real.
\end{itemize}

\subsection{Callbacks}
Callbacks implementados:
\begin{itemize}
  \item \textbf{ModelCheckpoint:} salva o melhor modelo conforme val\_accuracy;
  \item \textbf{EarlyStopping:} interrompe treinamento quando val\_loss para de melhorar;
  \item \textbf{ReduceLROnPlateau:} reduz learning rate quando val\_loss estagna.
\end{itemize}

\subsection{Mixed Precision (Opcional)}
Quando habilitado, mixed precision acelera treinamento em hardware compatível, reduzindo memoria e aumentando throughput. No projeto, o modo esta configurado como desativado por padrao.

\section{Avaliacao: Metrics e Analise de Erros}

\subsection{Accuracy e Loss}
Accuracy mede a proporcao de amostras corretamente classificadas. Loss quantifica o erro durante treinamento e avaliacao, sendo util para observar convergencia.

\subsection{Precision, Recall e F1-score}
Precision (para a classe prevista) e Recall (para a classe real) sao essenciais em cenarios desbalanceados. F1-score combina ambos em uma unica medida.

\subsection{Matriz de Confusao}
A matriz de confusao $C$ resume:
\begin{itemize}
  \item True Positives (acertos de Fall);
  \item True Negatives (acertos de Normal);
  \item False Positives (Normal previsto como Fall);
  \item False Negatives (Fall previsto como Normal).
\end{itemize}
Ela e a principal fonte visual de entendimento de erros.

\subsection{Leitura dos Graficos}
Os graficos de curves.png mostram:
\begin{itemize}
  \item train loss vs val loss;
  \item train accuracy vs val accuracy.
\end{itemize}
Quando val e train evoluem de forma semelhante, tende a indicar boa generalizacao (sem overfitting significativo).

\section{Inferencia em Tempo Real e Otimizacao}

\subsection{Buffer Circular de Sequencias}
Durante inferencia, o sistema acumula frames em um buffer ate completar T=20. A cada passo de inferencia, a janela atual e processada.

\subsection{Skip Frames}
O uso de inferencia a cada N frames (skip frames) reduz custo computacional e viabiliza execucao em CPU.

\subsection{Limiar e Confirmacao Temporal}
Para reduzir falsos positivos temporais, aplica-se:
\begin{itemize}
  \item um limiar mais conservador para considerar Fall;
  \item confirmacao temporal: so dispara Fall apos N predições consecutivas.
\end{itemize}

\subsection{TFLite e Quantizacao INT8}
O modelo e convertido para TFLite para desempenho. Quantizacao INT8 reduz tamanho e acelera em CPU, com possivel impacto na precisao (avaliar sempre).

\section{Visao Computacional com OpenCV}

\subsection{Captura de Video}
OpenCV utiliza VideoCapture para:
\begin{itemize}
  \item acessar webcam (indice 0);
  \item acessar arquivo de video (caminho do arquivo).
\end{itemize}

\subsection{Redimensionamento e Normalizacao}
A entrada do modelo precisa ser padronizada:
\begin{itemize}
  \item resize para HxW;
  \item conversao para float32 e normalizacao para [0,1];
  \item empilhamento temporal em formato correto para a rede.
\end{itemize}

\subsection{Visualizacao}
O sistema usa overlays com texto e paines para:
\begin{itemize}
  \item visualizar status do buffer;
  \item mostrar classe prevista e confianca;
  \item exibir alerta quando Fall for confirmado.
\end{itemize}

\section{Integração IoT: ESP32, Serial e MQTT}

\subsection{ESP32 e Arduino Framework}
O ESP32 executa firmware com:
\begin{itemize}
  \item controle de GPIO para buzzer e LEDs;
  \item recebimento de mensagem de alerta para ativar o alarme local.
\end{itemize}

\subsection{Comunicacao Serial (UART)}
Serial e recomendada para testes:
\begin{itemize}
  \item simples de configurar;
  \item nao exige WiFi;
  \item limita distancia e depende do cabo.
\end{itemize}
No contexto do projeto, usa-se configuracao de porta COM e baudrate.

\subsection{Comunicacao MQTT}
MQTT e um protocolo de mensageria publish/subscribe:
\begin{itemize}
  \item produtores publicam mensagens em um topico;
  \item assinantes recebem mensagens de um topico.
\end{itemize}
No projeto, o topico usado e ``fall\_detection/alerts''.

\subsection{Broker Mosquitto e WebSocket}
O Mosquitto e configurado com:
\begin{itemize}
  \item listener MQTT em 1883;
  \item listener WebSocket em 9001 para permitir conexao do cliente do app.
\end{itemize}

\subsection{Formato de Mensagem em JSON}
O payload do alerta e JSON. Conceitos importantes:
\begin{itemize}
  \item campos de controle: alert, confidence, timestamp;
  \item metadata opcional: frame\_id, model, etc.
\end{itemize}
Essa estrutura facilita interpretacao e depuracao, e permite evolucao do sistema sem quebrar compatibilidade.

\subsection{Estratégias de Robustez}
O estudo deve cobrir:
\begin{itemize}
  \item reconexao ao broker;
  \item tratamento de falhas (timeouts, portas incorretas, firewall);
  \item seguranca (quando usar MQTT com autenticação).
\end{itemize}

\section{App Mobile: React Native + Expo + MQTT over WebSocket}

\subsection{Por que WebSocket para MQTT no App}
Em ambiente JavaScript do React Native, conexoes TCP diretas nem sempre sao praticas. Assim, usa-se MQTT sobre WebSocket (WebSocket listener no broker) para permitir conectividade pelo ambiente mobile.

\subsection{Expo e Expo Go}
Expo simplifica:
\begin{itemize}
  \item configuracao do projeto;
  \item execucao no celular via Expo Go;
  \item acesso a API para audio, vibração e armazenamento.
\end{itemize}

\subsection{Gerenciamento de Estado com Context API}
O App utiliza Context API e useReducer para organizar:
\begin{itemize}
  \item estado de conexao MQTT;
  \item estado do alarme (ativo/inativo);
  \item eventos recebidos e persistidos.
\end{itemize}

\subsection{Servicos (Services)}
Arquitetura separa:
\begin{itemize}
  \item MqttService: conexao e subscribe no topico;
  \item AlarmService: som em loop e vibracao;
  \item EventStorage: persistencia no AsyncStorage.
\end{itemize}

\subsection{Fluxo do Alerta no App}
Fluxo conceitual:
\begin{enumerate}
  \item App recebe mensagem MQTT (JSON).
  \item Valida campo ``alert'' e ``confidence'' vs limiar.
  \item Aciona tela e inicia o alarme (som e vibração).
  \item Permite ao cuidador confirmar ou registrar falso alarme.
  \item Mantem historico persistido e filtro por tipo.
\end{enumerate}

\subsection{Historico e AsyncStorage}
AsyncStorage permite persistir eventos/configs localmente no dispositivo. Conceito central:
\begin{itemize}
  \item persistencia entre sessoes;
  \item consulta e filtros por tipo;
  \item limite de tamanho do historico (ex.: ate 200 eventos).
\end{itemize}

\section{Reprodutibilidade: Logs e Artefatos}

\subsection{Estrutura de Logs do Treinamento}
Cada treino gera uma pasta run-YYYYMMDD-HHMMSS/ contendo:
\begin{itemize}
  \item run\_metadata.json: hiperparametros e configuracoes;
  \item history.json / history.csv: evolucao por epoca;
  \item curves.png: graficos de perda e acuracia;
  \item classification\_report.txt: precision/recall/F1 por classe;
  \item confusion\_matrix.png: matriz visual;
  \item final\_metrics.json: loss e accuracy finais.
\end{itemize}

\subsection{Versionamento de Artefatos}
Boas praticas:
\begin{itemize}
  \item versionar artefatos leves (JSON/CSV/TXT/PNG);
  \item evitar versionar arquivos grandes (.h5 e dados brutos) ou usar Git LFS (se necessario).
\end{itemize}

\section{Limitações e Consideracoes de Ciencia do Mundo Real}

\subsection{Generalizacao}
Resultados excelentes no conjunto de teste podem nao refletir desempenho em:
\begin{itemize}
  \item outros ambientes/cameras;
  \item mudancas de iluminacao;
  \item posicoes e alturas distintas;
  \item variacoes de vestimenta e movimento.
\end{itemize}

\subsection{Falsos Positivos Temporais}
Deteccoes antecipadas podem ocorrer em transicoes de movimento. A mitigacao usa:
\begin{itemize}
  \item limiar de confidence mais alto;
  \item confirmacao por varias predições consecutivas.
\end{itemize}

\subsection{Performance}
Em hardware limitado:
\begin{itemize}
  \item usar TFLite/INT8;
  \item skip frames e buffer reduzido;
  \item evitar render pesado em loops (otimizar visualizacao).
\end{itemize}

\section{Referencias (para voce citar no Overleaf)}
Sugestoes de referencias utilizadas no contexto do projeto:
\begin{thebibliography}{9}
  \bibitem{kwolek2014} Kwolek, B.; Kepski, M. Human fall detection on embedded platform using depth maps and wireless accelerometer. Computer Methods and Programs in Biomedicine, 117(3), 489--501, 2014.
  \bibitem{mobnetv2} Sandler, M. et al. MobileNetV2: Inverted Residuals and Linear Bottlenecks. CVPR, 2018.
  \bibitem{lstm} Hochreiter, S.; Schmidhuber, J. Long Short-Term Memory. Neural Computation, 1997.
  \bibitem{lu2019} Lu, N. et al. Deep Learning for Fall Detection: Three-Dimensional CNN Combined With LSTM on Video Kinematic Data. IEEE JBHI, 2019.
  \bibitem{mubashir2013} Mubashir, M.; Shao, L.; Seed, L. A survey on fall detection: Principles and approaches. Neurocomputing, 2013.
  \bibitem{mqtt} MQTT.org. MQTT v3.1.1 Specification. (Documentacao oficial do protocolo).
  \bibitem{mosquitto} Eclipse Mosquitto. Documentacao e referencia do broker.
\end{thebibliography}

